\documentclass[11pt]{amsart}
\usepackage{amsmath,amssymb,amsthm,booktabs,hyperref,url}
\usepackage[margin=1.1in]{geometry}
\newtheorem{theorem}{Theorem}[section]
\newtheorem{proposition}[theorem]{Proposition}
\newtheorem{conjecture}[theorem]{Conjecture}
\newtheorem{lemma}[theorem]{Lemma}
\newtheorem{corollary}[theorem]{Corollary}
\theoremstyle{definition}
\newtheorem{observation}[theorem]{Observation}
\newtheorem{remark}[theorem]{Remark}
\newtheorem{example}[theorem]{Example}
\newcommand{\Q}{\mathbb{Q}}
\newcommand{\Z}{\mathbb{Z}}
\newcommand{\F}{\mathbb{F}}
\newcommand{\degphi}{\deg\varphi_E}
\newcommand{\ord}{\operatorname{ord}}
\newcommand{\Phig}{\Phi^{\mathrm{geom}}}

\title{Local component groups and the $\ell$-adic valuations of the modular degree}
\author{[Author]}
\address{[Affiliation]}
\email{[email]}
\date{\today}

\begin{document}

\begin{abstract}
Let $E/\Q$ be the $X_0(N)$-optimal curve of its isogeny class and $\varphi_E\colon X_0(N)\to E$ its modular parametrisation. For a bad prime $q$ let $\Phig_q$ denote the component group of the N\'eron model of $E$ over $\overline{\F}_q$. We study the inequality
\[
v_\ell(\degphi)\ \ge\ \sum_{q\mid N} v_\ell\bigl(\#\Phig_q\bigr)
\]
for odd primes $\ell$ at which $E[\ell]$ is irreducible. For $\ell\ge 5$ its multiplicative part is a consequence of a theorem of Kim and Ota on the Pollack--Weston conjecture combined with a result of Agashe, Ribet and Stein, under hypotheses that exclude $\ell=3$, primes $\ell\mid N$ and primes $q\equiv\pm1\pmod\ell$. We verify the inequality with no exception on every optimal curve of conductor at most $500000$ (more than two million curves) and on several hundred curves of conductor up to $4\cdot 10^6$ outside the tables, beyond all of these hypotheses. The new feature is the prime $\ell=3$: every prime of Kodaira type $\mathrm{IV}$ or $\mathrm{IV}^*$ contributes one to $v_3(\degphi)$, independently of its Tamagawa number. We prove the qualitative statement by a level-lowering argument, formulate a quantitative conjecture with a supplement for the ``twisted'' types $\mathrm{II}$, $\mathrm{II}^*$ and $\mathrm{I}_n^*$, and describe the exact failure of the inequality in the Eisenstein case.
\end{abstract}

\maketitle

\section{Introduction}

Let $E/\Q$ be an elliptic curve of conductor $N$ and let $f_E\in S_2(\Gamma_0(N))$ be the newform attached to it. Among the curves in the isogeny class of $E$ there is a distinguished one, the $X_0(N)$-optimal curve, for which the modular parametrisation $\varphi_E\colon X_0(N)\to E$ has minimal degree; throughout, $E$ denotes this curve and $\degphi$ its \emph{modular degree}. The modular degree is computable \cite{Cremona97,Watkins02} and its growth is understood (it is $N^{1+o(1)}$ times a power of the discriminant), but its prime factorisation is not. The purpose of this note is to describe, and partly prove, a local-to-global divisibility: the $\ell$-adic valuation of $\degphi$ is bounded below by local data at the bad primes.

For a bad prime $q$ let $\Phig_q = \Phi_q(\overline{\F}_q)$ be the group of components of the special fibre of the N\'eron model of $E$ at $q$, over the algebraic closure. Its order is read off the Kodaira type: $n$ for type $\mathrm{I}_n$; $1$ for $\mathrm{II},\mathrm{II}^*$; $2$ for $\mathrm{III},\mathrm{III}^*$; $3$ for $\mathrm{IV},\mathrm{IV}^*$; $4$ for $\mathrm{I}_0^*$ and $\mathrm{I}_n^*$ \cite[IV.9]{SilvermanATAEC}. The Tamagawa number $c_q=\#\Phi_q(\F_q)$ counts only the components defined over $\F_q$; the statements below are about $\Phig_q$, not $c_q$.

\begin{conjecture}\label{conj:main}
Let $E$ be $X_0(N)$-optimal and without complex multiplication, and let $\ell$ be an odd prime such that $E[\ell]$ is irreducible as a $\operatorname{Gal}(\overline{\Q}/\Q)$-module. Then
\begin{equation}\label{eq:main}
v_\ell(\degphi)\ \ge\ \sum_{q\mid N} v_\ell\bigl(\#\Phig_q\bigr).
\end{equation}
\end{conjecture}

For $\ell\ge5$ only multiplicative primes contribute, through $v_\ell(n_q)$ where $n_q=v_q(\Delta_E)$ is the exponent of the minimal discriminant. For $\ell=3$ the primes of type $\mathrm{IV}$ and $\mathrm{IV}^*$ contribute one each. The multiplicative part has a long history. Ribet and Takahashi \cite{RibetTakahashi97,Takahashi01} compared $\degphi$ with the degree of a parametrisation by a Shimura curve and found the product of the $n_q$ over the primes of the quaternion discriminant, up to primes where $E[\ell]$ is reducible; Pasten \cite{Pasten} controls that error term. Pollack and Weston \cite{PollackWeston11} conjectured, and Kim and Ota \cite{KimOta} proved, that moving a prime $q$ into the ``new'' part of the level changes the congruence ideal of $f_E$ by exactly the \emph{Tamagawa exponent} $t_q$, the largest $t$ with $\rho_{E,\ell}\bmod\ell^t$ unramified at $q$; for multiplicative $q$ one has $t_q=v_\ell(n_q)$ whether or not the reduction is split. Together with the theorem of Agashe, Ribet and Stein \cite{ARS12} that $\degphi$ divides the congruence number with quotient supported on primes whose square divides $4N$, this gives:

\begin{theorem}[Kim--Ota, Agashe--Ribet--Stein]\label{thm:known}
Let $\ell\ge5$ be a prime not dividing $N$ such that $\bar\rho_{E,\ell}$ is absolutely irreducible when restricted to $\operatorname{Gal}(\overline\Q/\Q(\sqrt{\ell^*}))$, $\ell^*=(-1)^{(\ell-1)/2}\ell$. Then
\[
v_\ell(\degphi)\ \ge \sum_{\substack{q\,\|\,N\\ q\not\equiv\pm1\ (\mathrm{mod}\ \ell)}} v_\ell(n_q).
\]
\end{theorem}

The data presented in Section~\ref{sec:data} say that none of the three restrictions in Theorem~\ref{thm:known} is needed for the inequality: it holds at $\ell=3$, at primes $\ell$ dividing $N$, and with the primes $q\equiv\pm1\pmod\ell$ included, in every one of more than two million optimal curves. The additive term at $\ell=3$ is of a different nature, since $\bar\rho_{E,3}$ is ramified at a prime of type $\mathrm{IV}$ or $\mathrm{IV}^*$ and no Tamagawa exponent is available. What makes it work is that the ramification is unipotent, so that the level can be lowered from $q^2$ to $q$:

\begin{theorem}\label{thm:IV}
Let $E/\Q$ have additive reduction of Kodaira type $\mathrm{IV}$ or $\mathrm{IV}^*$ at a prime $q\ge5$. Assume that $\bar\rho_{E,3}$ restricted to $\operatorname{Gal}(\overline\Q/\Q(\sqrt{-3}))$ is absolutely irreducible, and that $9\nmid N$. Then $3$ divides the modular degree of the $X_0(N)$-optimal curve in the isogeny class of $E$.
\end{theorem}

Theorem~\ref{thm:IV} is the qualitative half of the additive term; the quantitative half, namely that the contributions of several such primes and of the multiplicative primes add up, is Conjecture~\ref{conj:main}, and the data support it also at $q=3$ and when $9\mid N$.

Two further statements are supported by the data with no exception. The first concerns the \emph{twisted} types. Quadratic twisting at an odd prime $q$ exchanges $\mathrm{II}\leftrightarrow\mathrm{IV}^*$, $\mathrm{IV}\leftrightarrow\mathrm{II}^*$ and $\mathrm{I}_n\leftrightarrow\mathrm{I}_n^*$; at a prime of type $\mathrm{II}$, $\mathrm{II}^*$ or $\mathrm{I}_n^*$ the mod $3$ representation is a ramified quadratic character times a unipotent, that is, it has the shape of a twisted Steinberg representation. Give such a prime the weight $w_q=1$ for types $\mathrm{II},\mathrm{II}^*$ and $w_q=v_3(n_q)$ for type $\mathrm{I}_n^*$.

\begin{conjecture}\label{conj:twist}
With the hypotheses of Conjecture~\ref{conj:main} and $\ell=3$,
\[
v_3(\degphi)\ \ge\ \sum_{q\mid N} v_3\bigl(\#\Phig_q\bigr)\ +\ \Bigl(\sum_q w_q - \max_q w_q\Bigr).
\]
\end{conjecture}

In words, the twisted primes contribute all but the largest of their weights: two primes of type $\mathrm{II}$ force $3\mid\degphi$, one does not, three force $9\mid\degphi$. We have no mechanism for this statement.

The second statement concerns the Eisenstein case, in which Conjecture~\ref{conj:main} can fail.

\begin{conjecture}\label{conj:eis}
Let $E$ be optimal, not CM, with a rational $\ell$-isogeny, $\ell$ odd, and let $D$ be the deficit $\sum_q v_\ell(\#\Phig_q)-v_\ell(\degphi)$. Then
\[
D\ \le\ v_\ell\bigl(\#E(\Q)_{\mathrm{tors}}\bigr)+\bigl[\,E[3]\simeq\Z/3\oplus\mu_3\,\bigr],
\]
and also $D\le$ the number of curves in the isogeny class of $E$ with a rational point of order $\ell$. In particular $D\le0$ when $E$ has no rational $\ell$-torsion point, $D\le1$ for $\ell\ge5$, and $D\le2$ for $\ell=3$.
\end{conjecture}

Section~\ref{sec:local} collects the local facts, Section~\ref{sec:proof} proves Theorem~\ref{thm:IV}, Section~\ref{sec:data} describes the computations and Section~\ref{sec:two} discusses the prime $2$, where the analogous statements overlap with the literature on Watkins's conjecture.

\section{Local invariants}\label{sec:local}

Let $q$ be a prime of bad reduction. We use Tate's algorithm \cite[IV.9]{SilvermanATAEC} for the Kodaira type and $c_q$, and the following two facts.

\begin{lemma}\label{lem:tamagawa-exp}
Let $q$ be a prime of multiplicative reduction, $n_q=v_q(\Delta_E)$, and $\ell\ne q$ a prime. Then $\rho_{E,\ell}\bmod\ell^t$ is unramified at $q$ if and only if $\ell^t\mid n_q$. This holds for split and for nonsplit reduction.
\end{lemma}

\begin{proof}
Over the unramified quadratic extension of $\Q_q$ (which does not change the inertia group) $E$ is a Tate curve with parameter $u$, $v_q(u)=n_q$. The $\ell^t$-torsion is generated by $\mu_{\ell^t}$ and an $\ell^t$-th root of $u$, so inertia acts trivially on $E[\ell^t]$ exactly when $u$ is an $\ell^t$-th power in $\Q_q^{\mathrm{ur}}$ up to units, that is, when $\ell^t\mid n_q$.
\end{proof}

\begin{lemma}\label{lem:IV}
Let $q\ge5$ be a prime at which $E$ has reduction of type $\mathrm{IV}$ or $\mathrm{IV}^*$. Then the image of the inertia group $I_q$ in $\operatorname{Aut}(T_3E)$ is cyclic of order $3$, and its image in $\operatorname{Aut}(E[3])$ is a nontrivial unipotent subgroup. Consequently $E[3]^{I_q}$ is a line, $\bar\rho_{E,3}$ is tamely ramified at $q$, and the Serre conductor of $\bar\rho_{E,3}$ has exponent $1$ at $q$.
\end{lemma}

\begin{proof}
For $q\ge5$ the reduction is potentially good and the semistability defect, the order of the image of $I_q$ in $\operatorname{Aut}(T_\ell E)$ for any $\ell\ne q$, equals $12/\gcd(v_q(\Delta_E),12)$ \cite{Kraus90}; for $v_q(\Delta_E)\in\{4,8\}$ this is $3$. The image is tame since $q\ne3$, hence cyclic, generated by an element $\sigma$ of order $3$ with determinant $1$, so with eigenvalues the two primitive cube roots of unity. Hence $\sigma^2+\sigma+1$ annihilates $T=T_3E$, and $T$ is a module over $\Z_3[\sigma]/(\sigma^2+\sigma+1)\simeq\Z_3[\zeta_3]$, a discrete valuation ring with uniformiser $\pi=1-\zeta_3$ and residue field $\F_3$. Being torsion-free of $\Z_3$-rank $2$, $T$ is free of rank one over $\Z_3[\zeta_3]$, and $E[3]=T/3T\simeq\Z_3[\zeta_3]/(\pi^2)$ with $\sigma$ acting as multiplication by $\zeta_3=1-\pi$. Since $\pi\ne0$ in $\Z_3[\zeta_3]/(\pi^2)$, this is a nontrivial unipotent automorphism of order $3$; its fixed space is the line $\pi\Z_3[\zeta_3]/(\pi^2)$. The Serre conductor exponent at $q$ is $\dim E[3]-\dim E[3]^{I_q}=1$, there being no wild part.
\end{proof}

\begin{remark}
For the other additive types the image of inertia mod $3$ is not unipotent: it has order $2$ or $4$ (types $\mathrm{I}_0^*$, $\mathrm{III}$, $\mathrm{III}^*$), or it is minus a unipotent (types $\mathrm{II}$, $\mathrm{II}^*$, defect $6$) or a ramified quadratic character times a unipotent (type $\mathrm{I}_n^*$). The last two are the twisted Steinberg shapes of Conjecture~\ref{conj:twist}.
\end{remark}

\section{Proof of Theorem~\ref{thm:IV}}\label{sec:proof}

Let $f=f_E$, of level $N$ with $q^2\,\|\,N$, and $\bar\rho=\bar\rho_{E,3}$, which is irreducible by hypothesis and modular of weight $2$ and level $N$. By Lemma~\ref{lem:IV} the Serre conductor $N(\bar\rho)$ has exponent $1$ at $q$ and, as always, divides $N$ elsewhere, so $N(\bar\rho)\mid N/q$. By the level-lowering theorem of Ribet, Carayol and Diamond \cite{Ribet90,Carayol89,Diamond95}, applicable at $\ell=3$ under the hypothesis that $\bar\rho|_{\operatorname{Gal}(\overline\Q/\Q(\sqrt{-3}))}$ is absolutely irreducible, there is a newform $g$ of weight $2$ and level $M\mid N(\bar\rho)$, with coefficients in the ring of integers $\mathcal{O}$ of a number field, and a prime $\lambda\mid3$ of $\mathcal{O}$ such that $a_p(g)\equiv a_p(f)\pmod\lambda$ for all primes $p\nmid 3N$. Since $M<N$, the forms $g(dz)$, $d\mid N/M$, are old at level $N$ and lie in the orthogonal complement of $f$ in $S_2(\Gamma_0(N))$. A congruence between the Hecke eigenvalues of $f$ and of an eigenform in the orthogonal complement at almost all primes means that the Hecke algebra of level $N$ has a maximal ideal of residue characteristic $3$ containing the kernels of both eigencharacters, which is equivalent to $3$ dividing the congruence number $r_E$ of $f$ \cite[\S3]{ARS12}. Finally \cite[Theorem 2.1]{ARS12} gives $\degphi\mid r_E$ with every prime $p$ dividing $r_E/\degphi$ satisfying $p^2\mid 4N$; as $9\nmid N$, $v_3(\degphi)=v_3(r_E)\ge1$. \qed

\begin{remark}
The hypothesis $9\nmid N$ is used only in the last step and the data show the conclusion without it; so do they at $q=3$, where Lemma~\ref{lem:IV} does not apply because the ramification is wild. The proof gives one factor of $3$ for each such prime but says nothing about additivity with the multiplicative contributions or with each other; that is the content of Conjecture~\ref{conj:main} at $\ell=3$.
\end{remark}

\section{The computations}\label{sec:data}

\subsection{Data and protocol}
We used Cremona's tables \cite{Cremona97} in the form of the \texttt{ecdata} repository at commit \texttt{25cec5e}, which contain every elliptic curve over $\Q$ of conductor below $500000$ with its Kodaira symbols, Tamagawa numbers, modular degree, isogeny class, torsion and, up to conductor $400000$, the images of the mod $\ell$ Galois representations. Optimality is determined in the tables for every class of conductor at most $400000$ and for $359{,}009$ further classes up to $500000$. The statements were found on the curves of conductor at most $300000$, written down, and then tested once on the curves of conductor in $(300000,400000]$, which had been held out, and afterwards on the extension to $500000$. Nothing was adjusted after the hold-out test. The Kodaira types and discriminant exponents at the bad primes were recomputed with PARI/GP 2.15 \cite{PARI} from the $a$-invariants, and fifteen modular degrees spanning the configurations were recomputed with \texttt{ellmoddegree}, all agreeing with the tables.

\subsection{Results}
Table~\ref{tab:counts} gives, for each range and each small prime $\ell$, the number of optimal non-CM curves with $E[\ell]$ irreducible and a positive right-hand side in \eqref{eq:main}, the number of violations of Conjectures~\ref{conj:main} and~\ref{conj:twist}, and the number of Eisenstein curves with a positive bound together with the violations of Conjecture~\ref{conj:eis}. There are none.

\begin{table}[ht]
\centering\small
\begin{tabular}{llrrrrr}
\hline
set & $\ell$ & curves & $E[\ell]$ irreducible, bound $>0$ & violations C1 & violations C2 & Eisenstein, bound $>0$ / violations C3 \\
\hline
working set, $N \le 300000$ & 3 & 1,311,066 & 598,833 & 0 & 0 & 67,845 / 0 \\
 & 5 & 1,311,066 & 235,510 & 0 &  & 4,454 / 0 \\
 & 7 & 1,311,066 & 121,947 & 0 &  & 759 / 0 \\
 & 11 & 1,311,066 & 42,334 & 0 &  & 0 / 0 \\
 & 13 & 1,311,066 & 27,215 & 0 &  & 26 / 0 \\
\hline
hold-out, $300000 < N \le 400000$ & 3 & 429,936 & 206,749 & 0 & 0 & 18,225 / 0 \\
 & 5 & 429,936 & 79,550 & 0 &  & 989 / 0 \\
 & 7 & 429,936 & 41,387 & 0 &  & 137 / 0 \\
 & 11 & 429,936 & 14,244 & 0 &  & 0 / 0 \\
 & 13 & 429,936 & 9,212 & 0 &  & 2 / 0 \\
\hline
extension, $400000 < N \le 500000$ & 3 & 359,009 & 181,075 & 0 & 0 & 9,894 / 0 \\
 & 5 & 359,009 & 70,514 & 0 &  & 474 / 0 \\
 & 7 & 359,009 & 36,734 & 0 &  & 70 / 0 \\
 & 11 & 359,009 & 13,347 & 0 &  & 0 / 0 \\
 & 13 & 359,009 & 8,467 & 0 &  & 5 / 0 \\
\hline
out of table, $500000 < N \le 4\cdot 10^6$ & 3 & 154 & 110 & 0 & 0 & 0 / 0 \\
\hline
\end{tabular}

\caption{Verification counts. ``Bound $>0$'' means the right-hand side of \eqref{eq:main} is positive.}\label{tab:counts}
\end{table}

The theorem range of Theorem~\ref{thm:known} covers $281{,}739$ of the $913{,}244$ working-set pairs $(E,\ell)$ with $E[\ell]$ irreducible and a positive multiplicative bound; the remaining pairs have $\ell=3$ ($455{,}707$), $\ell\mid N$ ($175{,}798$ at $\ell\ge5$), or a contributing prime $q\equiv\pm1\pmod\ell$ ($410{,}074$), and show no violation either.

\subsection{Isolated configurations}
Table~\ref{tab:isolated} isolates the additive types at $\ell=3$: it lists, for curves with no multiplicative prime $q$ with $3\mid n_q$ and all additive primes of a single Kodaira type, the minimum of $v_3(\degphi)$ by type and multiplicity. The hold-out reproduces every minimum.

\begin{table}[ht]
\centering\small
\begin{tabular}{lccc}
\hline
type & one prime & two primes & three primes \\
\hline
IV & 1 (25,414) & 2 (969) & 3 (11) \\
IV$^*$ & 1 (30,168) & 2 (2,268) & 3 (45) \\
II & 0 (34,611) & 1 (2,073) & 2 (24) \\
II$^*$ & 0 (19,608) & 1 (968) & 2 (8) \\
III & 0 (34,904) & 0 (2,208) & 0 (34) \\
III$^*$ & 0 (25,574) & 0 (1,307) & 0 (29) \\
$\mathrm{I}_0^*$ & 0 (42,924) & 0 (3,093) & 0 (53) \\
$\mathrm{I}_n^*$ & 0 (133,322) & 0 (28,125) & 1 (1,483) \\
\hline
\end{tabular}

\caption{Minimum of $v_3(\degphi)$ (number of curves) over working-set curves whose additive primes are all of the given type, with no multiplicative $3$-contribution and $E[3]$ irreducible.}\label{tab:isolated}
\end{table}

The rows $\mathrm{IV}$, $\mathrm{IV}^*$ are Conjecture~\ref{conj:main}: each prime contributes exactly one at the minimum. The rows $\mathrm{II}$, $\mathrm{II}^*$ are Conjecture~\ref{conj:twist} with all weights equal to one. The row $\mathrm{I}_n^*$ mixes primes with and without $3\mid n$; separating them, two $\mathrm{I}_n^*$ primes with $3\mid n$ give $1$ and three give $2$. Across all configurations, the minimum of the excess over the bound of Conjecture~\ref{conj:main} as a function of (number of $\mathrm{II},\mathrm{II}^*$ primes, total $\mathrm{I}_n^*$ weight) is exactly $\sum w_q-\max w_q$ in every cell with at least two contributing primes, and the ``even subset'' rule of Ribet--Takahashi type (which would predict $2$ for two primes of weight $1$) is violated $8{,}979$ times.

\subsection{Sharpness and the rest of the degree}
The bounds are attained. Among working-set curves with $E[\ell]$ irreducible and a positive bound, $v_\ell(\degphi)$ equals the bound of Conjectures~\ref{conj:main} and~\ref{conj:twist} in $36.5\%$ of cases at $\ell=3$, $69.7\%$ at $\ell=5$ and $80.4\%$ at $\ell=7$, the proportion decreasing slowly with the conductor (from $75\%$ to $35\%$ at $\ell=3$ between conductors $300$ and $300000$). The whole odd part of $\degphi$ is another matter: for the $1{,}226{,}216$ working-set optimal non-CM curves with no isogeny of odd degree, the odd part of $\degphi$ equals $\prod_\ell \ell^{\,\text{bound}}$ in only $1.4\%$ of cases. The quotient is most often $9$, $3$, $15$, $45$, $27$ or $21$; its smallest prime factor is $3$ in two thirds of the cases and divides $N$ in about half. So the local component groups account for a definite, sharp part of the modular degree and for little of the rest, which is governed by congruences of other origins and by the size of the Petersson norm.

\subsection{Out of table}
To test the statements beyond the tables, we generated curves with prescribed local types and conductor between $500000$ and $4\cdot10^6$ (families with a type $\mathrm{IV}$ or $\mathrm{IV}^*$ prime, with two or three type $\mathrm{II}$ primes, and random curves), kept those with trivial isogeny class (so that the curve is optimal and $E[\ell]$ is irreducible for every $\ell$) and not CM, and computed their modular degrees with \texttt{ellmoddegree}. Table~\ref{tab:oot} summarises them; Conjectures~\ref{conj:main} and~\ref{conj:twist} hold for all of them.

\begin{table}[ht]
\centering\small
\begin{tabular}{lrrr}
\hline
family & curves & largest conductor & median seconds per degree \\
\hline
II,II & 86 & 3,968,800 & 0.4 \\
II,II,II & 1 & 2,832,200 & 0.2 \\
IV & 316 & 3,988,012 & 1.1 \\
IV* & 104 & 3,994,600 & 4.9 \\
IVmixed & 250 & 3,977,584 & 1.3 \\
random & 30 & 2,710,716 & 33.2 \\
all & 787 & 3,994,600 & 1.3 \\
\hline
\end{tabular}

\caption{Curves outside the tables, with modular degrees computed independently.}\label{tab:oot}
\end{table}

\subsection{The Eisenstein case}
Among the $75{,}430$ working-set optimal non-CM curves with a rational $3$-isogeny, $67{,}845$ have a positive bound and $2{,}777$ a positive deficit $D$; $D=2$ occurs $22$ times. No curve without a rational $3$-torsion point has a deficit, whatever its $3$-isogeny structure. With a rational point of order $3$ and mod $3$ image of Borel type, $D\le1$; $D=2$ occurs only when $E[3]\simeq\Z/3\oplus\mu_3$ (the split Cartan image \texttt{3Cs.1.1} in the notation of \cite{Sutherland16}, seventeen curves, all of rank $0$ and torsion $\Z/3$) or when $E$ has a point of order $9$ (five curves). The bounds of Conjecture~\ref{conj:eis} hold with no exception on all three ranges, and the torsion bound is attained with $D>0$ in $2{,}773$ cases. For $\ell=5$ and $7$ the deficits are at most $1$ and occur only with a rational $\ell$-torsion point ($192$ and $18$ cases in the working set).

\subsection{The twisted primes}
[[TWIST]]

\section{The prime 2}\label{sec:two}
At $\ell=2$ the modular degree is the subject of Watkins's conjecture $2^{\operatorname{rank}E(\Q)}\mid\degphi$ \cite{Watkins02} and of a substantial literature \cite{Dummigan06,CalegariEmerton09,Yazdani11,CaroPasten22,HatleyKundu,BKP}, and the mod $2$ representation has a rational $2$-isogeny for a large fraction of curves. We record only the following. For working-set optimal curves without a rational $2$-isogeny and not CM, the inequality $v_2(\degphi)\ge\sum_{q\,\|\,N}v_2(n_q)+\#\{q:\ \text{type }\mathrm{III},\mathrm{III}^*,\mathrm{I}_0^*\text{ or }\mathrm{I}_n^*\}$ holds with no exception on all three ranges ($964{,}432$, $325{,}510$ and $313{,}741$ curves), while counting $\mathrm{I}_0^*$ with weight $2=v_2(\#\Phig_q)$ fails ($123$, $26$ and $20$ exceptions). The isolated-configuration minima suggest that types $\mathrm{III}$, $\mathrm{III}^*$ and $\mathrm{I}_n^*$ contribute $3$ each and types $\mathrm{II}$, $\mathrm{IV}$ and $\mathrm{I}_0^*$ contribute $1$ each, in line with the theorem of Calegari and Emerton that a curve of odd modular degree has at most two odd bad primes and a rational $2$-torsion point or prime conductor \cite{CalegariEmerton09}. We have not attempted to reconcile the exponents with the results of \cite{Dummigan06,CaroPasten22,HatleyKundu,BKP}.

\section{Open questions}
\begin{enumerate}
\item Prove Conjecture~\ref{conj:main} at $\ell=3$ for multiplicative primes, and remove the conditions $\ell\nmid N$ and $q\not\equiv\pm1\pmod\ell$ from Theorem~\ref{thm:known}. The data say the inequality survives both; whether the equality of Kim and Ota does is a separate question.
\item Prove the quantitative additive statement: a Kim--Ota type formula in which the prime $q$ enters the level with exponent $2$ and the local invariant is the unipotent ramification of Lemma~\ref{lem:IV}.
\item Explain the ``all but the largest'' rule of Conjecture~\ref{conj:twist}. The quadratic twist by the character ramified at two twisted primes turns both into untwisted types, which suggests comparing the congruence modules of $f$ and of its twist, but the data exclude a congruence between $f$ and its own twist (Section~\ref{sec:data}).
\item Identify the deficit of Conjecture~\ref{conj:eis} with a natural invariant of the Eisenstein ideal at $\ell$, presumably through the cuspidal subgroup of $J_0(N)$ \cite{Mazur77}.
\end{enumerate}

\appendix
\section{Reproducibility}
All scripts, the pinned data commit and the result tables are in the repository \url{https://github.com/dsr001001/mcptrial} (directory \texttt{ecmine}). The checker takes a table of optimal curves with their Kodaira symbols, Tamagawa numbers, modular degrees, isogeny degrees and torsion, and reports the violations of each statement; it runs on the full tables in minutes. A single curve can be checked in PARI/GP with
\begin{verbatim}
E = ellinit([0,0,0,-99,360]); L = ellglobalred(E)[5]; d = ellmoddegree(E);
b = sum(i=1,#L, my(k=L[i][2]); if(k>=5, valuation(k-4,3), if(k==4||k==-4, 1, 0)));
print([valuation(d,3), b])
\end{verbatim}
where the codes $k=n+4$, $4$, $-4$ stand for $\mathrm{I}_n$, $\mathrm{IV}$, $\mathrm{IV}^*$. The pipeline was rerun from a clean clone and reproduced every count.

\section*{Acknowledgements}
[[The computations, the literature search and a first draft of this note were carried out with the assistance of an AI system (Claude, Anthropic); the author is responsible for all statements and proofs.]]

\bibliographystyle{amsplain}
\bibliography{refs}
\end{document}
